\documentclass{umk-matematika}
\begin{document}

\begin{center}
{\Large\bfseries Практическая работа 26}
\end{center}
\vspace{0.5cm}

\textbf{Задача 1.} Радиус шара $5$ см. Найдите площадь его поверхности.\par\vspace{10pt}
\textbf{Задача 2.} Радиус шара $3$ см. Найдите его объём.\par\vspace{10pt}
\textbf{Задача 3.} Площадь поверхности шара $64\pi$ см². Найдите радиус.\par\vspace{10pt}
\textbf{Задача 4.} Объём шара $288\pi$ см³. Найдите радиус.\par\vspace{10pt}
\textbf{Задача 5.} Сечение шара плоскостью, удалённой от центра на $8$ см, имеет площадь $36\pi$ см². Найдите радиус шара.\par\vspace{10pt}
\textbf{Задача 6.} Радиус шара $13$ см. Сечение находится на расстоянии $5$ см от центра. Найдите площадь сечения.\par\vspace{10pt}
\textbf{Задача 7.} Два шара имеют радиусы $3$ см и $6$ см. Во сколько раз объём большего шара больше объёма меньшего?\par\vspace{10pt}
\textbf{Задача 8.} Резервуар имеет форму шара радиусом $5$ м. Сколько краски потребуется для покраски его поверхности, если расход краски $0{,}2$ кг на 1 м²?\par\vspace{10pt}
\textbf{Задача 9.} Купол здания имеет форму полусферы радиусом $10$ м. Сколько краски потребуется для покраски внешней поверхности купола, если расход краски $0{,}2$ кг на 1 м²?\par\vspace{10pt}
\textbf{Задача 10.} Шар радиусом $15$ см переплавлен в цилиндр, радиус основания которого равен радиусу шара. Найдите высоту цилиндра.\par\vspace{10pt}
\newpage
\section*{Ответы}

\begin{enumerate}
  \item Ответ: $100\pi$ см²
  \item Ответ: $36\pi$ см³
  \item Ответ: $4$ см
  \item Ответ: $6$ см
  \item Ответ: $10$ см
  \item Ответ: $144\pi$ см²
  \item Ответ: в $8$ раз
  \item Ответ: $20\pi$ кг $\approx 62{,}8$ кг
  \item Ответ: $40\pi$ кг $\approx 125{,}6$ кг
  \item Ответ: $20$ см
\end{enumerate}
\end{document}
